% =====================================================================================================
% Related work (the \section and the label sec:related are in main.tex).
% Sources: bibliography/literature_review.md (verified state of the art) and the related-work parts of the
% theory notes (B3: Hofmann et al. paragraph and Rem. B3-conclusion; B2: Sec. B2-related; B1: Rem. B1-dragomir
% and Prop. B1-nofreelunch; A1: Sec. lb:statements; A3: Sec. lk-naive). Every key is in references.bib.
% Quotations checked in the downloaded texts: Hofmann et al. (arXiv 1506.03662, Sec. 3), Schmidt et al. (arXiv
% 1309.2388, Sec. 4.5), Zhou et al. (AISTATS 2019, experiments); Allen-Zhu (Sec. 7.1) as quoted in the A3 note.
% Numbers of the A4/A4b sweeps: final (results/A4_summary.json, results/A4b_summary.json, findings.md Sec. 7).
% A1 wording: assessment of hypothesis (A1) in theory/A1_lower_bound.tex (Sec. lb:statements; item (iii): the
% deterministic part of Thm. WS excludes, when kappa > 73, any gain of the type 1 + ln(n/kappa) for deterministic
% index sequences, cyclic ones included; randomly reshuffled methods are not settled). PEPit SAGA example:
% PEPit/examples/stochastic_and_randomized_convex_minimization/saga.py (gamma = 1/(2(mu n + L)), rate 1 - gamma mu).
% =====================================================================================================
\paragraph{Variance-reduced methods.}
In SAG~\cite{leroux2012stochastic,schmidt2017minimizing} the iterate follows a biased average of stored gradients,
in SVRG~\cite{johnson2013accelerating} stochastic gradients are corrected by a full gradient computed at a
periodically refreshed snapshot, and SAGA~\cite{defazio2014saga} stores one gradient per component and uses an
unbiased direction;
SDCA~\cite{shalevshwartz2013stochastic}, Finito and MISO~\cite{defazio2014finito,mairal2015incremental},
SARAH~\cite{nguyen2017sarah} and loopless versions of SVRG~\cite{hofmann2015variance,kovalev2020dont} complete the
family, surveyed in~\cite{gower2020variance}. These methods need $O((n+\kappa)\log(1/\varepsilon))$ component
gradients, with $\kappa$ defined from the largest smoothness constant of the components (Finito and MISO under a
big-data condition $n\gtrsim\kappa$). The proof of Proposition~\ref{prop:59} belongs to a family of analyses whose Lyapunov
function adds to $\norm{\theta-\etastar}^2$ a weighted sum of squared errors of the stored gradients, or of bounds
on them~\cite{schmidt2015nonuniform,hofmann2015variance,hu2017unified,gorbunov2020unified,kovalev2020dont},
whereas the original analysis of SAGA~\cite{defazio2014saga} uses Bregman divergences of the components; all certify
contraction factors of the form $1-\min\{\Theta(1/n),\Theta(1/\kappa)\}$, with constants that differ by small factors
between analyses. The analysis closest to the book is that of~\cite{hofmann2015variance}, which uses the same
square-norm inequality (``we sacrifice tightness'') and whose optimal step and rate are exactly the best ones that
the proof structure of Proposition~\ref{prop:59} can certify (\resref*{prop:B3-optimum}). Two sources of slack have
been identified: this square-norm inequality, and cocoercivity
with the worst-case constant $\max_iL_i$ (our $R^2$) instead of the average of the $L_i$, which non-uniform sampling
exploits~\cite{xiao2014proximal,schmidt2015nonuniform,needell2016stochastic}, or of an expected-smoothness
constant~\cite{gower2019sgd,gower2021stochastic}. Practical steps also exceed the theoretical ones: for SAG, the step
$1/L$ ``always converged and often performed better than the step size of $1/16L$ suggested by our
analysis''~\cite{schmidt2017minimizing}. We add a measurement along trajectories of the slack of every inequality of
the proof (the tightness ratios of \resref*{def:B3-ratios}), which also shows that strong convexity with the global
constant $\mu$ carries a sizeable share of the logarithmic gap when $\kappa\ge10n$ (Section~\ref{sec:exp-B3}).

\paragraph{Acceleration of variance-reduced methods.}
Indirect acceleration wraps a linearly convergent method in an inexact accelerated proximal-point loop:
Catalyst~\cite{lin2015universal,lin2018catalyst} and APPA~\cite{frostig2015unregularizing} turn the complexity of SAGA
or SVRG into $(n+\sqrt{n\kappa})\log(1/\varepsilon)$ up to logarithmic factors, and accelerated proximal
SDCA~\cite{shalevshwartz2014accelerated} has the same inner--outer structure. Direct optimal methods include
APCG~\cite{lin2015accelerated}, SPDC~\cite{zhang2015stochastic}, the primal--dual and gradient-extrapolation methods
of~\cite{lan2018optimal,lan2018random}, Varag~\cite{lan2019unified}, Point-SAGA~\cite{defazio2016simple}, which uses
proximal steps and no momentum, and Katyusha~\cite{allenzhu2018katyusha}. Katyusha writes Nesterov's method in the
linear-coupling form of~\cite{allenzhu2017linear} and adds to the coupling a pull towards the snapshot, the
``Katyusha momentum'', which keeps the query point close to the point where the estimator has no variance.
MiG~\cite{zhou2018simple} keeps only this negative momentum, SSNM~\cite{zhou2019direct} transfers it to SAGA, and
L-Katyusha~\cite{kovalev2020dont} replaces the outer loop by a coin flip and admits a one-step Lyapunov proof,
extended to arbitrary sampling in~\cite{qian2021lsvrg}; estimate sequences give the same
rate~\cite{kulunchakov2019estimate}. Our note on L-Katyusha (Appendix~\ref{sec:lkatyusha}) reproduces the proof
of~\cite{kovalev2020dont} with its constants (\resref*{thm:lk}); it adds the presentation in the style of the book,
the identification of the single exact cancellation performed by the negative momentum (\resref*{prop:cancel}), and
exact measurements of the certified contraction, of the variance bound and of the cancelled term along trajectories
(Section~\ref{sec:exp-A3}).

\paragraph{Why naive acceleration fails.}
Nesterov's momentum on the SVRG estimator is the accelerated proximal SVRG of~\cite{nitanda2014stochastic}, which
with one sample per iteration has the complexity of SVRG; growing mini-batches give a complexity of order
$(n+\min\{\kappa,n\sqrt\kappa\})\log(1/\varepsilon)$~\cite{nitanda2016accelerated}. In~\cite{allenzhu2018katyusha},
this scheme is identified with Katyusha without negative momentum and found, after careful tuning, to be sometimes
faster, but on some tasks ``without Katyusha momentum the algorithm does not even enjoy an accelerated convergence
rate''. Acceleration without a reference point is obtained in~\cite{driggs2022accelerating}, with mini-batches of
size of order $n^{2/3}$ and a dependence on $n$ of order $n^{2/3}$ instead of $\sqrt n$. The mechanisms invoked are error
accumulation in accelerated methods~\cite{schmidt2011convergence,devolder2014first}, an inner product in the tight
variance bound of the SVRG estimator, which the negative momentum cancels~\cite{allenzhu2018katyusha,zhou2019direct},
and the weakness of classical momentum with stochastic gradients~\cite{kidambi2018insufficiency}; the semidefinite
certificate of~\cite{hu2018dissipativity} for Katyusha is stated for negative-momentum weights of at least $1/5$.
To our knowledge, no lower bound excludes the acceleration of the naive scheme, and we do not claim one:
\resref{prop:naive-kappa} only shows that, without the cancellation, the Lyapunov argument of L-Katyusha certifies at
best a contraction $1-\kappa^{-2}$ per iteration. Our experiments agree
with~\cite{allenzhu2018katyusha}: divergence, or an exponent of $\kappa$ of about $0.8$ even with re-tuning at every
$\kappa$, on block quadratics, and competitive speed on logistic regression (Sections~\ref{sec:exp-A2}
and~\ref{sec:exp-A4}).

\paragraph{Lower bounds for finite sums.}
Lower bounds descend from the oracle model of~\cite{nemirovsky1983problem} and from Nesterov's chain quadratic, on
which a method whose iterates stay in the span of the observed gradients discovers one coordinate per
gradient~\cite[Theorem~2.1.13]{nesterov2018lectures} (a ``zero chain'' in the terminology of~\cite{carmon2020lower}).
For finite sums, deterministic incremental methods need $\Omega(n+\sqrt{n\kappa}\log(1/\varepsilon))$
gradients~\cite{agarwal2015lower}; span-based methods that sample the components i.i.d.\ from a fixed distribution
need $\Omega((n+\sqrt{n\kappa})\log(1/\varepsilon))$, in a dimension that grows with the number of
iterations~\cite{lan2018optimal}; the $\sqrt{n\kappa}\log(1/\varepsilon)$ term remains without the span assumption for
randomized algorithms with proximal oracles, for $\kappa$ larger than a constant multiple of $n$ (the
$n\log(1/\varepsilon)$ term becoming an additive $n$), and deterministic ones need $\Omega(n\sqrt\kappa\log(1/\varepsilon))$ at small enough
$\varepsilon$~\cite{woodworth2016tight}; and the $\sqrt{n\kappa}$ term holds, up to logarithmic factors, in every
dimension for oblivious methods, which include SAGA with a constant step~\cite{arjevani2016dimension}. The
$n\log(1/\varepsilon)$ term is proved for i.i.d.\ sampling: full-gradient steps can divide it by $1+\log(n/\kappa)$
when $\kappa\ll n$, and, for oblivious methods, not by more, up to logarithmic factors~\cite{hannah2018breaking}; the
corresponding bounds hold for loopless protocols such as those of L-SVRG and L-Katyusha in a dimension of order
$\sqrt{\kappa/n}\log(1/\varepsilon)$ only~\cite{han2024lower}. For
deterministic methods, cyclic ones included, the deterministic bound of~\cite{woodworth2016tight} already excludes a
gain of this type when $\kappa>73$; whether randomly reshuffled span-based methods, or deterministic ones when
$\kappa\le73$, can beat $n\log(1/\varepsilon)$ is not settled by these results. Our note
(Appendix~\ref{sec:lower_bound}) restates these results with their exact hypotheses in the notation of the book,
corrects hypothesis A1 accordingly, and builds an instance in the spirit of~\cite{agarwal2015lower,lan2018optimal}
(block-separable chains with a shared regularizer, so that every $\kappa>1$ is allowed) with the explicit dimension
condition of \resref{cor:count} and the explicit window of certified exponents of \resref{prop:A4}. Among the
works we reviewed, the only comment on the empirical exponent is in~\cite{zhou2019direct}: on two data sets,
multiplying $\kappa$ by $10$ slows accelerated methods down by about $3\approx\sqrt{10}$ and SAGA significantly more;
we fit the exponents over several decades of $\kappa$ on lower-bound instances (Section~\ref{sec:exp-A4}).

\paragraph{Tightness of convergence proofs and computer-assisted analyses.}
Performance estimation computes the worst case of a method by semidefinite programming~\cite{drori2014performance},
exactly for smooth strongly convex classes thanks to interpolation conditions~\cite{taylor2017smooth}; integral
quadratic constraints certify linear rates with small linear matrix inequalities~\cite{lessard2016analysis}; the best
rate certifiable by a quadratic Lyapunov function can be computed~\cite{taylor2018lyapunov}, also for stochastic
methods~\cite{taylor2019stochastic} and automatically~\cite{upadhyaya2024automated,upadhyaya2025autolyap}, and
PEPit~\cite{goujaud2024pepit} implements these tools. For variance reduction, the jump-system analysis
of~\cite{hu2017unified} certifies SAGA-type methods with a small linear matrix inequality whose Lyapunov function has
the structure of the one of Proposition~\ref{prop:59}, the dissipativity analysis of~\cite{hu2018dissipativity}
treats SVRG and Katyusha, and the SAGA example distributed with PEPit (file
\path{PEPit/examples/stochastic_and_randomized_convex_minimization/saga.py}, documented at
\url{https://pepit.readthedocs.io}) finds the one-step contraction proved in~\cite{defazio2014saga} tight for its
Bregman Lyapunov function at the
parameters it tests; the dual multipliers of such programs indicate which inequalities of a worst-case proof are
active~\cite{goujaud2023fundamental}. These tools compute the worst case over a class of functions, whereas our
tightness ratios measure how far each inequality of a given proof is from equality along one trajectory of one
instance. The two views are complementary: a guarantee can have the right order in the worst case
(\resref*{prop:B3-lower} shows that the exponent of $\kappa$ in Proposition~\ref{prop:59} cannot be improved for SAGA
with a constant step) and still be very pessimistic on the instances met in practice (Section~\ref{sec:exp-B3}). Our
optimization of $(\gamma,\Delta)$, \resref{prop:B3-optimum}, is the analytic counterpart of a Lyapunov search within
the two-parameter family of the book's proof; to our knowledge, the exact worst-case one-step contraction of the
Lyapunov function of Proposition~\ref{prop:59} has not been computed, and we do not compute it.

\paragraph{Local curvature and effective constants.}
The effective curvature $\muEff(\theta)$ of \resref{def:B2-eff} is the constant of the restricted secant
inequality~\cite{zhang2013gradient} at the point $\theta$, and linear convergence under such one-point or growth
conditions is classical~\cite{karimi2016linear,necoara2019linear}; the SAGA certificate of~\cite{hu2017unified} only
needs a sector condition at the minimizer, in line with our trajectory-dependent versions of
Proposition~\ref{prop:59}, the a-posteriori rates of \resref*{prop:B3-traj} and the a-priori rate on a region of
\resref*{prop:B3-region}. Our $\LEff$ is a
trajectory-wise version of the expected-smoothness constant~\cite{gower2019sgd,gower2021stochastic}, for which the
valid value given in~\cite{gower2021stochastic} for single-element uniform sampling is again $\max_iL_i$. For logistic
regression,
self-concordance compares Hessians at nearby points~\cite{bach2010self,sun2019generalized}, averaged stochastic
gradient descent adapts to the smallest eigenvalue of the Hessian at the optimum~\cite{bach2014adaptivity}, the
curvature at the optimum can be bounded through a measure of non-separability~\cite{freund2018condition}, SAGA-type
methods reach a local regime after identifying the active manifold of composite problems~\cite{poon2018local}, and
average-case analyses replace $\kappa$ by the spectral density of the Hessian~\cite{pedregosa2020acceleration}.
Faster-than-predicted convergence attributed to local curvature has been reported
repeatedly~\cite{lin2015universal,kulunchakov2019estimate,zhou2019direct}. Two results limit what effective constants
can give: one-point conditions alone do not allow acceleration~\cite{guilleescuret2022gradient}, and oblivious
methods need $\mu$ to accelerate~\cite{arjevani2017limitations}; our trajectory-dependent statements concern SAGA with
its fixed step. In the literature we reviewed we found no measurement of $\muEff$, or of the slack of each inequality
of a proof, along SAGA trajectories; Sections~\ref{sec:exp-B3} and~\ref{sec:exp-B2} provide them.

\paragraph{Preconditioning of stochastic methods.}
For a fixed symmetric positive definite matrix $P$, the step $\theta\mapsto\theta-\gamma P^{-1}g$ is a mirror step for
the reference function $\frac12\theta^\top P\theta$, governed by relative smoothness and relative strong
convexity~\cite{lu2018relatively,bauschke2017descent}; Bregman-SAGA with a quadratic reference function has a rate of
the form of Proposition~\ref{prop:59} with relative constants~\cite{dragomir2021fast} (compared with ours in
\resref*{rem:B1-dragomir}), and a change of variables transfers
Euclidean proofs to the fixed metric~\cite{kolte2015accelerating,gonen2016solving}. For ridge regression,
preconditioning by the regularized covariance reduces the relevant condition number to the dimension with sampling
proportional to the smoothness constants, or to $n$ times the largest ridge leverage score with uniform
sampling~\cite{gonen2016solving}. Preconditioned SVRG and Katyusha~X~\cite{liu2019acceleration} and the methods of
PROMISE~\cite{frangella2024promise} are analysed with relative or quadratic-regularity constants, the latter defined
relative to the Hessian at a reference point and close to one near the optimum. Quasi-Newton and diagonal
metrics~\cite{byrd2016stochastic,moritz2016linearly,sadiev2024stochastic} are analysed through eigenvalue bounds on the
metric only, so that their guarantees are worse than the Euclidean ones even when they help in practice, and
statistical preconditioning by subsampled Hessians has mostly been studied for distributed
optimization~\cite{shamir2014communication,hendrikx2020statistically}. In the worst case, second-order information on
the components does not remove the $\sqrt{n\kappa}$ term~\cite{arjevani2017oracle}, and relative smoothness alone does
not allow acceleration~\cite{dragomir2022optimal}. Our note (Appendix~\ref{sec:B1}) redoes the gradient-memory proof of
Proposition~\ref{prop:59} in the $P$-metric, which gives \resref{cor:B1-rate}, and shows in
\resref{prop:B1-nofreelunch} that, unlike for ridge regression, no fixed $P$ improves the global relative condition
number of logistic regression, so that the gains of preconditioning are entirely trajectory-dependent there.
